% Generated by Listas/make_statements.py from pratica-lecture-03.md, item 18. Do not edit.
\textbf{Source.} Mas-Colell, Whinston \& Green (1995), \emph{Microeconomic Theory}, Ch. 9 Exercises, Exercise 9.B.5\^{}B (E. Maskin), p.~302 {[}PDF 162{]}. Reproduced verbatim. Answer in the source: none in the book (see the MWG solutions manual in \texttt{Livros/}).



\begin{quote}

(E. Maskin) Consider a game with two players, player 1 and player 2, in which each player \(i\) can choose an action from a finite set \(M_i\) that contains \(m_i\) actions. Player \(i\)'s payoff if the action choices are \((m_1, m_2)\) is \(\phi_i(m_1, m_2)\).



\begin{enumerate}

\def\labelenumi{(\alph{enumi})}

\item

  Suppose, first, that the two players move simultaneously. How many strategies does each player have?

\item

  Now suppose that player 1 moves first and that player 2 observes player 1's move before choosing her move. How many strategies does each player have?

\item

  Suppose that the game in (b) has multiple SPNEs. Show that if this is the case, then there exist two pairs of moves \((m_1, m_2)\) and \((m_1', m_2')\) (where either \(m_1 \neq m_1'\) or \(m_2 \neq m_2'\)) such that either

\end{enumerate}



\[\text{(i)}\quad \phi_1(m_1, m_2) = \phi_1(m_1', m_2')\]



or



\[\text{(ii)}\quad \phi_2(m_1, m_2) = \phi_2(m_1', m_2').\]



\begin{enumerate}

\def\labelenumi{(\alph{enumi})}

\setcounter{enumi}{3}

\item

  Suppose that for any two pairs of moves \((m_1, m_2)\) and \((m_1', m_2')\) such that \(m_1 \neq m_1'\) or \(m_2 \neq m_2'\), condition (ii) is violated (i.e., player 2 is never indifferent between pairs of moves). Suppose also that there exists a pure strategy Nash equilibrium in the game in (a) in which \(\pi_1\) is player 1's payoff. Show that in any SPNE of the game in (b), player 1's payoff is at least \(\pi_1\). Would this conclusion necessarily hold for any \emph{Nash} equilibrium of the game in (b)?

\item

  Show by example that the conclusion in (d) may fail either if condition (ii) holds for some strategy pairs \((m_1, m_2)\), \((m_1', m_2')\) with \(m_1 \neq m_1'\) or \(m_2 \neq m_2'\) or if we replace the phrase \emph{pure strategy Nash equilibrium} with the phrase \emph{mixed strategy Nash equilibrium}.

\end{enumerate}

\end{quote}



\textbf{Optional extras.}

- Tadelis 7.5 (veto power, p.~148), 7.7 (roommates voting, p.~149) and 7.8 with 8.3 (Brothers, pp.~149, 170): more trees, strategy sets and SPE.

- Tadelis 7.3 (tic-tac-toe, pp.~147--148): counting information sets in a large tree; needs a spreadsheet.

- Tadelis 8.1 and 8.2 (p.~170): mixed-strategy SPE of the voluntary Battle of the Sexes (read Fig. 8.2, p.~154) and of mutually assured destruction (read §8.3.3, pp.~163--166, Fig. 8.11); 8.1 is the slides' Example 2 of Lecture 4.

- Muñoz-Garcia Ch. 4 Ex. 2 (entry-predation, p.~156), Ex. 5--6 (sequential Pareto coordination and chicken, pp.~166, 169), Ex. 10 (trust and reciprocity, p.~177), Ex. 19--20 (backward induction drills, pp.~202, 204): all with worked answers.

- Gibbons 2.2 (Samaritan's dilemma, pp.~130--131) and 2.4 (two-period contribution game, pp.~131--132): continuous-action backward induction after item 17.



\begin{center}\rule{0.5\linewidth}{0.5pt}\end{center}

